\chapter{Architectural Approach}

\section{Chosen Framework}

The selected architectural design framework is \textbf{Attribute-Driven Design (ADD)}.

ADD fits this project because the VerdeMart scenario is mainly driven by measurable quality
attribute pressures. The main question is not only which components should exist, but how the
commerce core should behave when surrounding operational systems are delayed, stale,
unavailable, or recovering.

In this project, ADD is applied as follows:

\begin{enumerate}
    \item Drivers identified: availability, consistency, performance, recoverability, and modifiability.
    \item Six quality attribute scenarios defined, each with a measurable response.
    \item Tactics selected per scenario: asynchronous integration, transactional outbox-style event publishing, read-model separation, retry queue, reconciliation, adapter pattern, and observability.
    \item Target architecture decomposed around those tactics: nopCommerce Core, Integration Layer, Operational Adapters, External Systems, and Support and Operations View.
    \item Validation focused on the riskiest path: asynchronous propagation, degraded dependency handling, retry, and recovery.
\end{enumerate}

\section{Justification}

The chosen scenario requires nopCommerce to evolve from a web storefront into the commerce
core of a wider operational ecosystem. This creates architectural pressure in five main areas:
availability, consistency, performance, recoverability, and modifiability.

ADD is appropriate because each of these pressures can be expressed as a concrete quality
attribute scenario. Warehouse failure drives the need for asynchronous fulfillment and retry.
Stale inventory drives the need for inventory visibility, revalidation, and reconciliation. Order
bursts drive the need to protect customer-facing paths from slow backend synchronization.
Shipping outage drives the need for durable retry and explicit recovery state. New operational
channels drive the need for adapters and stable integration boundaries.

This gives a direct trace from the problem to the architecture:

\begin{center}
\begin{tabular}{|p{0.2\textwidth}|p{0.2\textwidth}|p{0.45\textwidth}|}
\hline
\textbf{ADD input} & \textbf{Scenario} & \textbf{Architectural consequence} \\
\hline
Availability & QAS 1 & Avoid synchronous dependency on warehouse systems during checkout. \\
\hline
Consistency & QAS 2, QAS 6 & Introduce explicit inventory state, stale-state visibility, conflict detection, and reconciliation. \\
\hline
Performance & QAS 3 & Keep checkout and order reads separate from slow integration processing. \\
\hline
Recoverability & QAS 4 & Store failed integration work durably and retry after dependency recovery. \\
\hline
Modifiability & QAS 5 & Add external systems through adapters instead of direct database coupling. \\
\hline
\end{tabular}
\end{center}

\section{Why Not ACDM or ADM as the Main Framework}

\textbf{ACDM} was considered because it is useful for architectural decision-making, risk
discussion, and stakeholder validation. However, it is not the main framework for this project
because the strongest design input is the set of measurable quality attribute scenarios. The
project must show runtime behavior under failure, delay, stale state, and recovery; ADD gives a
more direct path from those scenarios to concrete architectural tactics and component
responsibilities.

ACDM-style reasoning is still useful later in the project, especially when documenting ADRs and
explaining rejected alternatives. For example, the decision to use asynchronous integration can
be recorded as an ADR with alternatives such as direct synchronous calls or shared database
access. However, the decision itself is primarily driven by the availability and recoverability
scenarios, which makes ADD the better primary method.

\textbf{ADM/TOGAF} was also considered, but it is broader than necessary for this assignment.
ADM is better suited to enterprise-wide transformation across business, data, application, and
technology architecture. VerdeMart does have an omnichannel business context, but the scope of
this work is a focused software architecture evolution of nopCommerce, not a full enterprise
architecture programme.

\section{Tactics Selected per Quality Attribute}

ADD requires explicit tactics to address each quality attribute scenario. The following table
records which tactics were chosen and why each was preferred over the alternatives
considered in the ADRs.

\begingroup
\small
\setlength{\tabcolsep}{4pt}
\renewcommand{\arraystretch}{1.15}
\begin{center}
\begin{tabular}{|p{0.19\textwidth}|p{0.30\textwidth}|p{0.40\textwidth}|}
\hline
\textbf{Quality attribute} & \textbf{Tactics selected} & \textbf{Architectural manifestation} \\
\hline
Availability & Asynchronous integration, durable recording of work & Outbox pattern (ADR 2): order and its integration task are committed atomically; warehouse unavailability does not block checkout. \\
\hline
Consistency & Stale-state marking, explicit reconciliation & Inventory metadata fields (\texttt{IsStale}, \texttt{ConflictFlag}); reconciliation record created on divergence; checkout revalidates against \texttt{PendingReconciliation}. \\
\hline
Performance & Separate synchronous and asynchronous paths, read-model projection & Customer request path contains no external I/O; integration work runs outside the HTTP pipeline in the Integration Worker. \\
\hline
Recoverability & Idempotency, retry with backoff, dead-letter handling & Each integration task carries a correlation identifier and idempotency key; the retry scheduler applies exponential backoff; failed work beyond the retry limit is moved to a dead-letter queue visible to operators. \\
\hline
Modifiability & Adapter pattern, stable event contracts, no shared database & Each external system is reached through a dedicated adapter (ADR 3); internal message contracts are owned by the Integration Coordination context; no adapter accesses nopCommerce tables directly. \\
\hline
\end{tabular}
\end{center}
\endgroup

\section{How ADD Shapes the Target Architecture}

Using ADD means that each architectural element in the target design should answer a specific
quality attribute pressure. The surrounding systems are not added only because they exist in the
business scenario; they are added because they create design pressure that nopCommerce must
handle explicitly.

The target architecture emphasises asynchronous integration because QAS 1 requires that a
warehouse timeout does not block checkout. A direct synchronous call from checkout to the
warehouse cannot satisfy that constraint, so accepted orders are propagated through a durable
integration path.

Status is made explicit for pending, degraded, failed, and recovered work because QAS 1
and 4 require support agents and operators to understand what is happening during warehouse or
shipping failure. Retry and reconciliation are part of the architecture because QAS 2, QAS 4,
and 6 assume stale inventory, contradictory cross-channel stock, failed label creation, and
later recovery. External systems are connected through adapters because QAS 5 requires new
operational channels to be added without rewriting checkout or sharing the nopCommerce database.

The next chapter uses these ADD-driven tactics to define the target architecture and its
component responsibilities.
